\section{CLAUDE.md 与项目配置}
%% ============================================================

CLAUDE.md 是项目级的持久化配置，相当于给 Claude Code 的"项目说明书"。

\subsection{用 /init 生成然后精简（Tip \#28）}

CLAUDE.md 是项目根目录下的 Markdown 文件，给 Claude 提供持久指令：构建命令、编码规范、架构决策、仓库约定。每次会话开始时 Claude 都会读取它。

运行 \texttt{/init} 自动生成初始版本，但输出通常很臃肿。精简原则：如果你说不出某一行为什么存在，就删掉它。

\subsection{每一行都用 Litmus Test 检验（Tip \#29）}

\begin{importantbox}{CLAUDE.md 的黄金法则}
对 CLAUDE.md 的每一行都问：没有这行，Claude 会犯错吗？如果 Claude 本来就做得对，这行就是噪声。大约有 150--200 条指令的"预算"，超过后遵从度下降，而系统 prompt 已经占用了约 50 条。
\end{importantbox}

\subsection{让 Claude 自己更新 CLAUDE.md（Tip \#30）}

当 Claude 犯错时，说"update the CLAUDE.md file so this doesn't happen again"。Claude 会写入自己的规则，下次自动遵守。随着时间推移，CLAUDE.md 变成一个由真实错误驱动的活文档。

为防止文件无限膨胀，用 \texttt{@imports} 引用外部文件（如 \texttt{@docs/solutions.md}），CLAUDE.md 本身保持精简。

\subsection{用 .claude/rules/ 做条件加载（Tip \#31）}

在 \texttt{.claude/rules/} 目录下放置 Markdown 文件，按主题组织规则。通过 frontmatter 中的 \texttt{paths} 字段控制加载时机：

\begin{lstlisting}[caption={条件加载规则示例}, language=yaml]
---
paths:
  - "**/*.ts"
  - "**/*.tsx"
---
# TypeScript conventions
- Use strict mode
- Prefer interfaces over types
\end{lstlisting}

TypeScript 规则只在 Claude 处理 \texttt{.ts} 文件时加载，Go 规则只在处理 \texttt{.go} 文件时加载。

\subsection{用 @imports 按需引入（Tip \#32）}

用 \texttt{@docs/git-instructions.md}、\texttt{@README.md}、\texttt{@package.json} 甚至 \texttt{@\textasciitilde/.claude/my-project-instructions.md} 引用外部文档。Claude 在需要时才读取，不会每次会话都加载。

\subsection{本章小结}

CLAUDE.md 的核心理念：精简、针对性、活文档。用 \texttt{/init} 起步后立刻精简；每行都做 litmus test；让 Claude 从错误中自学；用 \texttt{.claude/rules/} 和 \texttt{@imports} 分散配置，保持主文件简洁。

%% ============================================================
